\section{Conclusion}
\label{sec:conclusion}

This paper proposes frequent losers---firms that never win yet
participate in abnormally many public tenders---as a screening marker
for cover bidding in procurement auctions. Using 4.5 million
tender-items from Brazil's BEC platform (2009--2019), we find strong
support for FL as a screening marker (H1): FL-present tenders exhibit
a robust conditional price premium of 3.6--7.7\% across OLS,
cross-fitting, and matching estimators, and FL firms co-participate
with CADE-convicted cartelists at 3.5 times the expected rate
($p < 0.001$). Multiple diagnostics provide suggestive evidence for
cover bidding as the underlying mechanism (H2): the FL--price
effect concentrates in competitive markets
(Proposition~\ref{prop:market_selection}), Bajari--Ye tests reject
exchangeability and conditional independence of FL bid residuals, and
structural estimation favors Regime~2
($\hat{\sigma}_c / \hat{\sigma}_g = 0.72$). No single test is
conclusive, but the consistency across independent approaches
supports the cover-bidding interpretation.

\paragraph{Identification and caveats.}
Our estimates capture a conditional association, not a causal effect.
The OLS, cross-fit, and matching approaches all rely on selection on
observables (plus fixed effects); unobserved market-level factors
that attract both FL firms and higher prices cannot be fully excluded.
Sensitivity analysis ($RV = 17.5\%$) bounds the scope for
omitted-variable bias, and the cross-fit estimate (3.6\%) removes any
mechanical link between the FL classification rule and the price
outcome, but neither eliminates the possibility of confounding.
The staggered DiD is underpowered (MDE = 0.076); the FL definition
necessarily involves misclassification; and strategic adaptation
(rotating cover bidders, allowing occasional wins) is a concern for
operational deployment.

The network-split results reveal that concentrated-market FL firms
show no significant price effect ($-0.018$, $p > 0.3$), while
competitive-market FL firms drive the entire effect (0.126,
$p < 0.01$). Three interpretations merit consideration: (i)~cartels
in concentrated markets sustain prices through market power rather
than cover bidding; (ii)~some FL firms in concentrated markets are
genuinely unsuccessful entrants rather than cover bidders; or
(iii)~item fixed effects absorb the effect in thin markets with few
observations per cell. Distinguishing among these requires richer
data on cartel structure. For screening purposes, this heterogeneity
implies that the FL marker is most informative in competitive
markets---precisely where detection is most valuable.

\paragraph{Magnitude in context.}
The 3.6--7.7\% range falls at the low end of the cartel overcharge
distribution documented by \citet{connor2007price} (median 23\%
across 674 cartels). This is consistent with FL presence marking
cartel-affected tenders rather than measuring the full overcharge:
fixed effects absorb within-item price variation unrelated to FL, and
the cross-fit estimate further discounts any in-sample mechanical
correlation. The true cartel markup likely exceeds the FL coefficient.
As an illustrative exercise, applying the OLS coefficient to FL-present
procurement spending yields an estimated welfare loss of R\$734M
(2.4\% of BEC spending) over 2009--2019; the cross-fit estimate
implies R\$400M (1.3\%). These figures are upper bounds conditional
on the causal interpretation and should be treated as order-of-magnitude
benchmarks, not precise policy estimates
(Table~\ref{tab:welfare_bounds}, Appendix~\ref{app:extensions}).

\paragraph{Policy implications.}
The FL screen requires only participation and outcome data, produces
firm-level flags for targeted investigation, and complements bid-level
screens in a two-stage workflow. FL flags are a starting point for
investigation, not standalone enforcement evidence.

The screen complements rather than substitutes for leniency programs.
Leniency targets insiders willing to defect; FL screening targets
the observable footprint of cover-bidding operations from the outside.
In jurisdictions where leniency uptake is low---as in much of Latin
America---proactive screening may be the only feasible detection
channel. Conversely, FL flags can inform leniency outreach by
identifying firms most likely involved in active cartels.

\paragraph{Portability.}
The FL screen's data requirements---participation records and
outcomes---are available in most electronic procurement systems.
The approach is portable to any setting where cover bidding
generates repeat-losing firms: other Brazilian states, EU member
states with centralized e-procurement (TED), and developing economies
adopting digital procurement platforms. The key institutional
requirement is that participation records include firm identifiers
and outcome flags; bid-level data are not needed.

\paragraph{Future research.}
Four directions merit investigation: external validation beyond BEC
(other states, other countries, federal procurement);
deeper network analysis (shared representatives, financial links);
dynamic models of FL-screen--leniency interaction; and field
experiments in which procurement agencies adopt FL screening.
